В работе была рассмотрена генерация юмора как открытая творческая задача, для которой трудно построить внешний детерминированный верификатор. Центральная идея состояла в том, что юмор не следует сводить к имитации одного эталонного ответа. Для одного и того же запроса может существовать множество удачных шуток, поэтому более естественной является многореференсная постановка, в которой модель обучается увеличивать вероятность наблюдаемого множества эталонных ответов.

Для этой цели был разработан и реализован полный программный конвейер. Он включает сбор и очистку корпуса шуток, удаление дублей и почти совпадающих текстов, построение векторных представлений, извлечение ключевых слов, поиск эталонных ответов, генерацию кандидатов, обучение LoRA-адаптеров и автоматическую попарную оценку. Исходная коллекция содержала приблизительно \(664\,000\) шуток; после дедупликации итоговый корпус составил приблизительно \(570\,000\) шуток. На его основе было построено около \(14\,600\) многореференсных запросов.

На методологическом уровне работа расширяет идею обучения без внешнего верификатора на случай нескольких эталонных ответов. В идеальной формулировке награда задаётся вероятностной массой, которую модель назначает всему наблюдаемому множеству эталонных шуток. Поскольку полные шутки являются длинными последовательностями, сырые вероятности таких последовательностей оказываются численно неудобными и смещёнными в пользу коротких ответов. Поэтому в практической реализации использовалась нормированная логарифмическая масса эталонных ответов, вычисляемая через принудительное учительское декодирование.

Экспериментальная часть проверила предложенный подход на моделях Qwen3-1.7B и Qwen3-4B. Для обеих моделей обучались LoRA-адаптеры, а качество оценивалось через автоматические попарные сравнения с использованием Gemini 2.5 Flash в роли модели-оценщика. Для каждой размерности сравнивались четыре варианта: базовая модель без рассуждения, базовая модель с рассуждением, MRVF-модель без рассуждения и MRVF-модель с рассуждением.

Результаты оказались различными для двух размеров модели. Для Qwen3-1.7B MRVF не дал уверенного улучшения относительно лучшего базового режима: первое место заняла базовая модель с рассуждением, а MRVF-модель с рассуждением заняла второе место с близким результатом. Этот результат показывает, что предложенный метод не является автоматическим улучшением для любой базовой модели и чувствителен к масштабу модели, качеству рассуждений и бюджету обучения.

Для Qwen3-4B результат оказался положительным. Обе MRVF-модели превзошли обе базовые модели, а MRVF-модель с рассуждением заняла первое место с долей побед 0.523. Это поддерживает основную гипотезу работы: многореференсная логарифмическая цель может служить полезным обучающим сигналом для генерации юмора без отдельной модели награды. Особенно важно, что базовый режим рассуждения для Qwen3-4B сам по себе оказался хуже прямой генерации, но после MRVF-обучения вариант с рассуждением стал лучшим. Это указывает на то, что полезным является не сам факт длинного рассуждения, а обучение распределения рассуждений, связанных с удачными финальными шутками.

В то же время результаты следует интерпретировать осторожно. Эффект на Qwen3-4B положительный, но умеренный: преимущество в прямых попарных сравнениях с базовыми моделями составляет несколько процентных пунктов. Оценка проводилась автоматически, а не через человеческую разметку. Кроме того, наблюдаемые множества эталонных шуток неполны и не исчерпывают все возможные хорошие ответы. Следовательно, MRVF оптимизирует суррогатную цель, связанную с корпусными эталонами, но не является прямой оптимизацией человеческой смешности.

Работа также выявила несколько практических ограничений. Во-первых, качество результатов существенно зависит от шаблона генерации: модель должна возвращать именно короткую финальную шутку, а не объяснение или повторяющийся текст. Во-вторых, в финальных обучающих запусках рассуждения почти всегда достигали заданного максимального бюджета, что указывает на необходимость дальнейшего контроля длины и структуры рассуждений. В-третьих, использовались короткие отладочные запуски, размер группы \(G=2\) и два эталонных ответа на запрос; более надёжная проверка требует более крупных запусков и абляций.

Основной вклад работы состоит в том, что она переносит идею обучения без внешнего верификатора с задач с коротким проверяемым ответом на открытую творческую задачу с несколькими допустимыми ответами. В работе показано, как построить многореференсный набор данных из неразмеченного корпуса шуток, как сформулировать устойчивую логарифмическую цель для эталонных ответов и как проверить её влияние на реальные генерации модели. Положительный результат на Qwen3-4B показывает, что это направление заслуживает дальнейшего исследования.

Дальнейшая работа должна идти в нескольких направлениях. Во-первых, необходимо провести более длинные обучающие запуски и абляции по числу траекторий, числу эталонных ответов, коэффициенту эталонного слагаемого и бюджету рассуждения. Во-вторых, требуется человеческая оценка хотя бы на подмножестве примеров, чтобы проверить согласованность автоматической оценки с восприятием людей. В-третьих, полезно улучшить построение эталонных множеств: учитывать веса поиска, фильтровать слабые эталоны и строить более разнообразные наборы шуток. Наконец, следует сравнить MRVF не только с базовыми моделями, но и с дообучением с учителем и другими методами постобучения.